\section{Conversion to simple unweighted graphs}\label{sec:unweighted}

We include an explicit deterministic conversion so that rational edge
weights and loops introduce no additional complexity assumption. The
comparison uses the original, unnormalized weight scale.
We replace each vertex by a cluster and approximate each weighted edge
by a bipartite graph whose edge count on every rectangle approximates the
weighted-density prediction with a uniform absolute error.
An arbitrary cut may split each cluster. Its cluster proportions then
describe a random cut of the original graph, so the rectangle estimates
control all cuts of the enlarged graph, not just cuts constant on clusters.

\begin{lemma}\label{lem:unweighted}
Let $G$ be a graph on $N$ vertices with nonnegative rational weights
$w_{uv}$ on unordered distinct vertex pairs, and
$\sum_{u<v}w_{uv}\le1$. For every fixed $\epsilon'>0$, one can construct
in deterministic polynomial time a simple unweighted graph $G'$ and an
integer scale $Q$ such that
\[
 \left|\frac{\operatorname{MaxCut}(G')}{Q}
             -\operatorname{MaxCut}(G)\right|\le\epsilon'.
\]
The same conclusion holds when the input graph has loops, after deleting
them without renormalizing its other weights.
\end{lemma}

\begin{proof}
Choose a prime $p>\max\{2,2N^2/\epsilon'\}$. A prime of polynomial size can be found deterministically
by testing successive integers, using Bertrand's postulate to bound the
search. Replace every vertex by a cluster indexed by $\mathbb F_p^6$.
For each $u<v$ let $I_{uv}\subseteq\mathbb F_p$ be any fixed set of
$\lfloor p w_{uv}\rfloor$ residues, for example the first such residues
in the usual ordering. Join $x$ in cluster $u$ to $y$ in cluster $v$
exactly when
\[
 x\cdot y\in I_{uv}.
\]
There are no edges within clusters. This defines a simple graph on
$Np^6$ vertices, constructible by enumerating all possible pairs.
Take $Q=p^{12}$.

We first estimate its edge count on arbitrary rectangles between two
clusters. Fix a residue set $I$ and subsets $S,T\subseteq\mathbb F_p^6$.
Writing $\psi(a)=\exp(2\pi i a/p)$, Fourier inversion gives
\[
 \one_I(a)=\sum_{h\in\mathbb F_p}\beta_h\psi(ha),\qquad
 \beta_h=\frac1p\sum_{a\in I}\psi(-ha),\qquad
 |\beta_h|\le1.
\]
For $h\ne0$, the matrix $M_h(x,y)=\psi(hx\cdot y)$ satisfies
\[
 M_hM_h^*=p^6 I,
\]
by character orthogonality, so its operator norm is $p^3$. Consequently,
\begin{align*}
 \left|\#\{(x,y)\in S\times T:x\cdot y\in I\}
                -\frac{|I|}{p}|S||T|\right|
 &\le\sum_{h\ne0}|\beta_h|p^3\sqrt{|S||T|}\\
 &\le p^{10}.
\end{align*}
For $I=I_{uv}$, division by $p^{12}$ and
$\bigl||I|/p-w_{uv}\bigr|\le1/p$ show that
\begin{equation}\label{eq:rectangle-discrepancy}
 \left|\frac{e_{uv}(S,T)}{p^{12}}
       -w_{uv}\frac{|S|}{p^6}\frac{|T|}{p^6}\right|
 \le\frac1p+\frac1{p^2}.
\end{equation}

Now fix an arbitrary cut of $G'$, and let $a_u\in[0,1]$ be the
fraction of cluster $u$ on its positive side. Its crossing edges between
clusters $u,v$ form two rectangles. Equation~\eqref{eq:rectangle-discrepancy}
therefore compares their scaled count with
\[
 w_{uv}\bigl(a_u(1-a_v)+(1-a_u)a_v\bigr)
\]
to error at most $2/p+2/p^2$.
Summing over the at most $N(N-1)/2$ pairs costs at most
$N(N-1)(1/p+1/p^2)$. The displayed weighted expression, summed over
pairs, is the expected cut weight obtained by independently assigning
vertex $u$ a positive sign with probability $a_u$. It is at most
$\operatorname{MaxCut}(G)$. This proves the desired upper bound for
every cut of $G'$.

Conversely, give all vertices in each cluster the sign of a maximum cut
of $G$. The same rectangle bound compares its scaled weight with
$\operatorname{MaxCut}(G)$ to at most the same error, giving the reverse
inequality. Our choice of $p$ ensures
$N(N-1)(1/p+1/p^2)<\epsilon'$, which proves the assertion.
Loops contribute zero to every cut and can be deleted at the outset.
\end{proof}
